% 05_wps.tex - Implementacion: WPs, tareas, entregables, hitos y Gantt.
\section{Plan de implementacion}
% SSOT del plan: data/workpackages.yaml (lo edita research-architect con la
% skill grant-workpackages). Esta seccion DERIVA del YAML: prohibido cambiar
% estructura o cifras aqui sin actualizar antes el YAML.

\subsection{Paquetes de trabajo}
% TODO: por cada WP del YAML: titulo, objetivos (IDs O1..On de
% context/hypothesis_matrix.md), tareas y entregables con su criterio de
% aceptacion. Los meses y plazos de los hitos CASAN con el calendario de la
% convocatoria (context/call_guidelines.md): si la convocatoria no define
% fecha, \todo{DATO PENDIENTE: [plazo de la convocatoria]}.

\subsection{Entregables y hitos}
% TODO: tabla resumen de entregables (id, descripcion, WP, criterio de
% aceptacion) y de hitos (id, descripcion, dependencias, mes). Con booktabs
% exclusivamente (skill latex-tables): sin lineas verticales ni \hline.

\subsection{Gantt del proyecto}
% Gantt derivado de data/workpackages.yaml (convention latex-tables: booktabs
% puro, sin pgfgantt). Tn = tramos del calendario (trimestres o meses, segun
% la duracion de la convocatoria): ajustar el numero de columnas a la
% duracion real y marcar con X las celdas activas de cada WP. Duracion y
% meses: SIEMPRE de context/call_guidelines.md; si no constan, marcador.
\begin{table}[htbp]
\centering
\caption{Gantt de los paquetes de trabajo (derivado de data/workpackages.yaml).}
\label{tab:gantt}
\begin{tabular}{llccccc}
\toprule
WP & Titulo & T1 & T2 & T3 & T4 & T5 \\
\midrule
WP1 & TODO: titulo del WP1 & \todo{DATO PENDIENTE: meses de WP1} & & & & \\
WP2 & TODO: titulo del WP2 & & & & & \\
WP3 & TODO: titulo del WP3 & & & & & \\
\bottomrule
\end{tabular}
\end{table}

\subsection{Gestion del proyecto}
% TODO: organizacion (comite de direccion, coordinacion del consorcio,
% seguimiento de entregables) y gestion de contingencias. En una PPC:
% mecanismos de coordinacion publico-privada y toma de decisiones.
